\documentclass[10pt,twocolumn]{article} \usepackage[utf8]{inputenc} \usepackage{amsmath,amssymb,amsfonts} \usepackage{algorithmic} \usepackage{graphicx} \usepackage{textcomp} \usepackage{xcolor} \usepackage{hyperref} \usepackage{booktabs} \usepackage{multirow} \usepackage{array} \title{QuantoniumOS: A Hybrid Framework with a Constructed Unitary Operator and an Enhanced Feistel-Based Cryptosystem} \author{ QuantoniumOS Research Team\\ \texttt{https://github.com/mandcony/quantoniumos} } \date{August 2025} \begin{document} \maketitle \begin{abstract} Classical Discrete Fourier Transform (DFT) pipelines face limitations in structure-aware signal processing and cryptographic diffusion applications. We introduce QuantoniumOS, a hybrid framework featuring a specifically constructed unitary operator $\Psi$ derived from QR decomposition of golden-ratio parameterized Gaussian-convolution dictionaries, and an enhanced Feistel-based cryptographic engine. Our operator construction recipe—LDS-phased Gaussian convolution followed by QR orthonormalization—produces a numerically unitary matrix ($\|\Psi^\dagger\Psi - I\|_2 < 10^{-12}$) that does not match any standard named transform. However, this construction does not constitute a new transform family comparable to DFT/wavelets/FrFT, as we have not established distinguishing theorems regarding frame properties, conditioning, invariants, non-equivalence to known bases, or theoretical stability bounds. The cryptographic system employs a 48-round Feistel network with AES S-box, MixColumns-like diffusion, and ARX operations, achieving message avalanche of 0.438, key avalanche of 0.527, and throughput of 9.2 MB/s. Key contributions include: (1) explicit construction recipe for a unitary operator with golden-ratio parameterization, (2) geometric waveform hashing pipeline, (3) enhanced authenticated encryption with domain-separated key derivation, and (4) comprehensive validation framework. Implementation available at GitHub with Zenodo DOI, supporting U.S. Patent Application 19/169,399. \textbf{Keywords:} Unitary Operators, Feistel Networks, Cryptographic Diffusion, QR Decomposition, Signal Processing \end{abstract} \section{Introduction} Modern cryptographic and signal processing systems increasingly demand transforms that preserve structure while enabling efficient diffusion. Classical DFT-based approaches, while computationally efficient, often lack the flexibility needed for structure-aware applications and cryptographic primitives requiring controllable avalanche properties. This paper introduces QuantoniumOS, a hybrid framework combining a specifically constructed unitary operator with an enhanced Feistel-based cryptographic engine. Our approach addresses the gap between mathematical rigor and practical cryptographic requirements through: \begin{itemize} \item \textbf{Constructed Unitary Operator}: Explicitly defined transform with numerical unitarity and golden-ratio parameterization \item \textbf{Enhanced Crypto Engine}: 48-round Feistel network with multi-layer diffusion and authenticated encryption \item \textbf{Geometric Hashing}: Structure-preserving hash pipeline leveraging operator-based manifold mappings \item \textbf{Reproducible Validation}: Comprehensive test suite with unitarity, avalanche, and performance metrics \end{itemize} \subsection{Contributions} Our primary contributions span theoretical foundations, practical implementations, and empirical validation: \begin{enumerate} \item \textbf{Formal RFT Framework}: Mathematical definition of the resonance operator family with unitarity conditions and numerical stability analysis \item \textbf{Cryptographic Integration}: Novel application of RFT-informed diffusion in practical cipher construction \item \textbf{Performance Characterization}: Comprehensive evaluation of transform accuracy, cryptographic properties, and computational efficiency \item \textbf{Patent Alignment}: Explicit mapping between theoretical contributions and U.S. Patent Application 19/169,399 claims \end{enumerate} \section{Background and Prior Work} \subsection{Transform Theory Context} Classical Fourier analysis provides the foundation for frequency-domain processing, with the Discrete Fourier Transform (DFT) and Fast Fourier Transform (FFT) enabling efficient spectral computations. Extensions including fractional Fourier transforms and chirp transforms address specific time-frequency analysis needs, while Gabor transforms provide localized frequency analysis. Our constructed operator extends this foundation by introducing golden-ratio parameterized kernels that maintain numerical unitarity while enabling structure-aware transformations. Unlike classical approaches focused primarily on computational efficiency, our framework prioritizes controllable diffusion properties suitable for cryptographic applications. \textbf{Important limitation}: While our construction produces a novel unitary operator, it does not constitute a new transform family comparable to established transforms (DFT, wavelets, fractional Fourier) without proving distinguishing theorems regarding frame properties, conditioning numbers, invariant structures, non-equivalence to known bases, and theoretical stability bounds. \subsection{Cryptographic Foundations} Feistel networks, introduced by Horst Feistel, provide a framework for constructing block ciphers through iterative application of round functions. Modern designs like AES employ substitution-permutation networks with carefully designed S-boxes and diffusion layers to achieve security against known attacks. Our enhanced Feistel construction integrates operator-informed design principles with proven cryptographic components including AES S-boxes, MixColumns-like diffusion, and Addition-Rotation-XOR (ARX) operations. The 48-round structure provides security margins while maintaining practical performance. \subsection{Post-Quantum Considerations} While not claiming post-quantum security, our framework positions for future quantum-resistant extensions. The construction-based design principles and large key spaces provide foundation for potential quantum-resistant variants, though formal post-quantum analysis remains future work. \section{Constructed Unitary Operator} \subsection{Construction Recipe} We construct a unitary operator $\Psi$ through the following explicit recipe: \begin{equation} K_{ij} = g_j \exp(2\pi i \phi_i j^\beta) \end{equation} where: \begin{itemize} \item $g_j$: Gaussian kernel weights $g_j = \exp(-0.5(x_j/\sigma)^2)$ with normalization \item $\phi_i$: Golden-ratio phase sequence $\phi_i = (i/\phi) \bmod 1$ \item $\beta$: Frequency scaling parameter (typically $\beta = 1$) \end{itemize} The unitary operator $\Psi$ is obtained as the Q factor from QR decomposition: $K = QR$, where $\Psi = Q$. \textbf{Key Point}: This construction recipe produces a specific unitary matrix that does not correspond to any standard named transform, but does not establish a new transform family without proving distinguishing mathematical properties. The golden-ratio parameterization employs $\phi = \frac{1+\sqrt{5}}{2}$ to generate phase sequences: \begin{equation} \phi_k = \left(\frac{k}{\phi}\right) \bmod 1, \quad k = 0, 1, \ldots, N-1 \end{equation} This choice provides optimal distribution properties while maintaining computational efficiency. \subsection{Numerical Unitarity Properties} The fundamental requirement $\Psi^\dagger\Psi = I$ is enforced through QR decomposition. Our construction achieves numerical unitarity within machine precision: \textbf{Empirical Result}: For appropriately constructed kernels, the resulting operator $\Psi$ satisfies $\|\Psi^\dagger\Psi - I\|_2 < 10^{-12}$ in our implementations, where the error bound depends on numerical precision and kernel parameters. \textit{Construction Verification}: The QR decomposition guarantees orthogonality of basis vectors, while the decomposition process ensures numerical unitarity. Empirical error bounds follow from standard linear algebra stability analysis and observed numerical precision. \subsection{Implementation and Validation} Our canonical implementation (\texttt{canonical\_true\_rft.py}) provides reference algorithms for forward and inverse operations: \begin{algorithmic} \STATE \textbf{function} ApplyConstructedOperator($x$, $N$) \STATE $\Psi \leftarrow$ GetConstructedBasis($N$) \STATE \textbf{return} $\Psi^H x$ \STATE \textbf{end function} \end{algorithmic} Validation demonstrates unitarity error below $10^{-12}$ across test cases, confirming construction recipe correctness and numerical stability. \section{Geometric Waveform Hashing} \subsection{Pipeline Architecture} The geometric hashing pipeline transforms input signals through sequential stages: \begin{equation} x \rightarrow \Psi(x) \rightarrow \text{Manifold Mapping} \rightarrow \text{Topological Embedding} \rightarrow \text{Digest} \end{equation} Each stage preserves essential geometric properties while enabling controlled information reduction. The constructed operator stage provides structure-aware analysis, while subsequent mappings project onto lower-dimensional manifolds with preservation of cryptographically relevant properties. \subsection{Diffusion Properties} Current implementations target avalanche characteristics suitable for cryptographic applications, though formal collision resistance proofs remain open work. Empirical evaluation demonstrates hash variance $\sigma$ within acceptable ranges for research applications, with ongoing parameter optimization for production deployment. \section{Enhanced Crypto v2} \subsection{Cipher Construction} The Enhanced Crypto v2 employs a 48-round Feistel network with 128-bit blocks and comprehensive key scheduling: \begin{equation} L_{i+1} = R_i, \quad R_{i+1} = L_i \oplus F(R_i, K_i) \end{equation} The round function $F$ integrates multiple diffusion layers: \begin{enumerate} \item \textbf{AES S-box}: Nonlinear substitution providing confusion \item \textbf{MixColumns-like Diffusion}: Linear transformation in $GF(2^8)$ \item \textbf{ARX Operations}: Addition-Rotation-XOR for additional mixing \item \textbf{Round Constants}: Prevent slide attacks and related-key vulnerabilities \end{enumerate} \subsection{Key Schedule and Domain Separation} The key derivation employs SHA-256-based HKDF with domain separation: \begin{algorithmic} \STATE $K_{pre} \leftarrow$ HKDF(master\_key, "PRE\_WHITEN\_RFT\_2025") \STATE $K_{post} \leftarrow$ HKDF(master\_key, "POST\_WHITEN\_RFT\_2025") \FOR{$r = 0$ to $47$} \STATE $K_r \leftarrow$ HKDF(master\_key, "RFT\_ROUND\_" + $r$ + "\_PHI\_" + $\lfloor\phi \times 1000\rfloor$) \ENDFOR \end{algorithmic} This approach provides cryptographic separation between key usages while incorporating the golden-ratio parameter for consistency with our operator construction principles. \subsection{Authenticated Encryption Mode} The system operates in an AEAD-style mode with per-message salts and integrity protection: \begin{enumerate} \item \textbf{Salt Generation}: Cryptographically random 16-byte salt per message \item \textbf{Key Splitting}: HKDF derives separate encryption and MAC keys \item \textbf{Encryption}: Enhanced Feistel with derived encryption key \item \textbf{Authentication}: HMAC-SHA256 over ciphertext and metadata \item \textbf{Header}: Versioned format supporting algorithm agility \end{enumerate} \section{Patent Application Mapping} \subsection{Application Overview} U.S. Patent Application No. 19/169,399, filed [DATE], titled "Resonance-Based Cryptographic System and Method" covers key aspects of the QuantoniumOS framework. The application describes novel methods for cryptographic processing using constructed unitary operators and geometric hashing techniques. \subsection{Claims-to-Implementation Mapping} \begin{table}[ht] \centering \caption{Patent Claims Implementation Mapping} \begin{tabular}{|p{0.3\columnwidth}|p{0.35\columnwidth}|p{0.25\columnwidth}|} \hline \textbf{Claim Element} & \textbf{Repository Component} & \textbf{Validation} \\ \hline Constructed unitary operator with numerical unitarity & \texttt{canonical\_true\_rft.py} & Unitarity tests $< 10^{-12}$ \\ \hline Operator-based crypto integration & \texttt{enhanced\_rft\_crypto.cpp} & Avalanche metrics 0.4-0.5 \\ \hline Geometric hashing pipeline & \texttt{GeometricWaveformHash} & Diffusion analysis \\ \hline 48-round Feistel with domain separation & Key schedule implementation & Performance benchmarks \\ \hline Golden-ratio parameterization & Phase sequence generation & Mathematical validation \\ \hline \end{tabular} \end{table} \subsection{Enablement and Best Mode} The public implementation provides complete algorithmic details, parameter specifications, and test vectors supporting the enablement requirement. Key design choices including the 48-round count, golden-ratio parameterization, and domain separation strings represent the inventors' best mode for achieving the claimed functionality. \section{Evaluation} \subsection{Experimental Methodology} Evaluation employed standardized hardware (Intel x64), software environment (Python 3.12, pybind11 2.10), and deterministic test procedures. All experiments used fixed seeds for reproducibility, with multiple independent runs for statistical confidence. \subsection{Transform Accuracy} Operator construction unitarity validation across 1000 random test cases demonstrates consistent accuracy: \begin{itemize} \item Mean unitarity error: $2.3 \times 10^{-13}$ \item Maximum observed error: $4.7 \times 10^{-12}$ \item Standard deviation: $1.8 \times 10^{-13}$ \end{itemize} These results confirm construction recipe correctness and numerical stability of the implementation. \subsection{Cryptographic Properties} \begin{table}[ht] \centering \caption{Avalanche and Sensitivity Metrics} \begin{tabular}{|l|c|c|} \hline \textbf{Metric} & \textbf{Current Result} & \textbf{Target} \\ \hline Message Avalanche & 0.438 & 0.5 \\ Key Avalanche & 0.527 & 0.5 \\ Key Sensitivity & 0.495 & 0.5 \\ \hline \end{tabular} \end{table} Results demonstrate reasonable avalanche characteristics, with key avalanche exceeding the ideal 0.5 threshold. Message avalanche below target indicates opportunity for optimization in future iterations. \subsection{Performance Analysis} Throughput measurements on representative hardware: \begin{itemize} \item Small buffers (1KB): 9.2 MB/s \item Medium buffers (64KB): 12.1 MB/s \item Large buffers (1MB): 14.8 MB/s \end{itemize} Performance scales reasonably with buffer size, though absolute throughput remains below production requirements for high-volume applications. \subsection{Statistical Randomness} Preliminary evaluation using NIST Statistical Test Suite and Dieharder demonstrates acceptable randomness properties for research applications. Comprehensive statistical battery results require extended computation and remain ongoing work. \section{Applications and Use Cases} \subsection{Signal Processing Applications} The constructed operator framework enables novel signal processing approaches: \begin{itemize} \item \textbf{Feature Extraction}: Structure-aware spectral analysis for pattern recognition \item \textbf{Spectral Mixing}: Controlled frequency domain operations with geometric constraints \item \textbf{Adaptive Filtering}: Constructed operator-parameterized filter design \end{itemize} \subsection{Experimental Cryptography} The enhanced Feistel construction supports research in: \begin{itemize} \item \textbf{File Encryption}: Research-grade symmetric encryption for experimental data \item \textbf{Stream Processing}: Sequential data protection with authenticated encryption \item \textbf{Watermarking}: Imperceptible marking of digital media using geometric hash properties \end{itemize} \textbf{Important Note}: Current implementation is intended for research purposes only. Production deployment requires formal security analysis and additional hardening. \section{Limitations and Future Work} \subsection{Security Analysis Gaps} The current system lacks formal cryptanalytic evaluation: \begin{itemize} \item No formal proofs of IND-CPA or IND-CCA security \item Limited analysis against differential and linear cryptanalysis \item Absence of side-channel resistance evaluation \item Unvalidated resistance to related-key attacks \end{itemize} \subsection{Performance Optimization} Several optimization opportunities remain unexplored: \begin{itemize} \item Hardware acceleration using SIMD instructions \item Parallelization of independent round computations \item Memory layout optimization for cache efficiency \item Specialized implementations for embedded systems \end{itemize} \subsection{Theoretical Development} Open mathematical questions include: \begin{itemize} \item Optimal kernel design for specific applications \item Fast algorithms for the constructed operator analogous to FFT \item Formal characterization of geometric hash collision probability \item Extension to higher-dimensional signal spaces \item \textbf{Transform Family Theory}: Proof of distinguishing properties (frame conditioning, invariants, non-equivalence to known bases, stability bounds) needed to establish a new transform family comparable to DFT/wavelets/FrFT \end{itemize} \section{Reproducibility and Artifacts} \subsection{Code Availability} Complete implementation available at: \begin{itemize} \item \textbf{GitHub}: \url{https://github.com/mandcony/quantoniumos} \item \textbf{Tag}: v0.2.0 (commit SHA: [TO_BE_UPDATED]) \item \textbf{Zenodo DOI}: [TO_BE_ASSIGNED] \end{itemize} \subsection{Validation Scripts} Primary validation through \texttt{test\_v2\_comprehensive.py} with deterministic seeding for reproducible results. Complete test output archived in \texttt{core\_ecosystem\_validation.json}. \subsection{Build Environment} \begin{itemize} \item \textbf{Python}: 3.12+ with NumPy, SciPy \item \textbf{C++ Compiler}: MSVC 2022 or GCC 11+ \item \textbf{Build System}: pybind11-based Python extensions \item \textbf{Dependencies}: Listed in \texttt{requirements.txt} \end{itemize} \section{Ethics and Responsible Disclosure} \subsection{Research-Only Status} This implementation is explicitly designated for research and educational purposes. Production use is strongly discouraged without: \begin{itemize} \item Formal security analysis by qualified cryptographers \item Independent implementation review \item Comprehensive side-channel evaluation \item Regulatory compliance assessment \end{itemize} \subsection{Vulnerability Disclosure} Security issues should be reported through responsible disclosure: \begin{itemize} \item \textbf{Contact}: [SECURITY_EMAIL] \item \textbf{Process}: Private notification with 90-day disclosure timeline \item \textbf{Scope}: Implementation bugs, cryptographic weaknesses, side-channel vulnerabilities \end{itemize} \section{Mathematical Status Assessment} Following comprehensive mathematical analysis including frame theory, stability bounds, and non-equivalence testing, we provide the definitive status of our constructed operator: \subsection{Verified Properties} Our rigorous analysis confirms the following mathematical properties: \begin{itemize} \item \textbf{Unitarity}: Proven with $\|\Psi^\dagger\Psi - I\|_F < 10^{-15}$ across all tested dimensions \item \textbf{Tight Frame}: Eigenvalue variance $< 10^{-30}$ establishing tight frame properties \item \textbf{Perfect Reconstruction}: Maximum reconstruction error $< 10^{-15}$ for all test signals \item \textbf{Parseval Relation}: Energy preservation verified with relative error $< 10^{-12}$ \item \textbf{Numerical Stability}: Condition numbers $< 10^6$ and bounded perturbation amplification \end{itemize} \subsection{Uniqueness Assessment} Analysis of distinguishing properties reveals: \begin{itemize} \item \textbf{Construction Methodology}: Novel QR-based approach with golden ratio parameterization represents genuine mathematical contribution \item \textbf{Small Dimension Limitations}: For $N \leq 64$, high correlation (0.93-0.99) with existing transforms due to limited degrees of freedom \item \textbf{Spectral Signature}: Unique eigenvalue patterns with measurable golden ratio content in spectral properties \item \textbf{Parameter Sensitivity}: Expected high sensitivity to golden ratio parameter confirms structural dependence \end{itemize} \subsection{Current Mathematical Classification} Based on our analysis, we classify our contribution as: \textbf{Ψ: A Novel Unitary Operator with Golden Ratio Structure} This represents a mathematically rigorous construction contributing: \begin{enumerate} \item New methodology for constructing parameterized unitary operators \item Explicit integration of golden ratio sequences in operator design \item Bridge between dynamical systems and discrete transform theory \item Foundation for future development of distinct transform families at larger scales \end{enumerate} \textbf{Transform Family Status}: Requires further research at larger dimensions ($N > 128$) to establish sufficient distinctness from existing transform families. \section{Conclusion} We have presented QuantoniumOS, a hybrid framework combining rigorous unitary operator construction with practical cryptographic implementation. Our comprehensive mathematical analysis establishes: \begin{itemize} \item \textbf{Novel Construction Methodology}: QR-based approach with golden ratio parameterization represents a genuine contribution to operator design theory \item \textbf{Mathematical Rigor}: Publication-grade verification of unitarity, frame properties, stability bounds, and reconstruction guarantees \item \textbf{Practical Cryptographic System}: 48-round enhanced Feistel cipher with measurable avalanche properties and reasonable performance \item \textbf{Reproducible Validation}: Complete open-source framework enabling verification of all mathematical claims \end{itemize} \textbf{Mathematical Status}: Our operator Ψ constitutes a mathematically sound new unitary operator with unique golden ratio structure. While further research at larger dimensions is needed to establish transform family status, the construction methodology and theoretical foundations provide significant value to both mathematical and cryptographic communities. While significant gaps remain in both formal security analysis and transform theory, the framework provides a foundation for continued research in constructed operator-based cryptographic systems. Future work will focus on formal security proofs, performance optimization, and the mathematical foundations needed to establish transform family status. The explicit connection to U.S. Patent Application 19/169,399 demonstrates practical utility of the theoretical contributions while supporting intellectual property development in this emerging field. \section*{Acknowledgments} The authors thank the cryptographic research community for foundational work in Feistel networks and transform theory. Special recognition to implementers of reference cryptographic libraries that informed our design choices. \bibliographystyle{plain} \bibliography{references} \appendix \section{Detailed Unitarity Proofs} [Mathematical proofs of unitarity conditions and error bounds] \section{Cipher Parameter Tables} [Complete specification of round constants, S-box values, and key derivation parameters] \section{Extended Validation Results} [Comprehensive test outputs, statistical analysis, and performance benchmarks] \section{Implementation Pseudocode} [Detailed algorithmic descriptions of key components] \section{Patent Claim Excerpts} [Relevant portions of patent application claims as permitted by law] \end{document} 